% !TeX root = ../main.tex
\section{Phụ lục: Hướng dẫn cài đặt và chạy hệ thống}\label{sec:phuluc}

Phụ lục này mô tả cách dựng lại toàn bộ hệ thống trên một máy mới, từ cơ sở dữ liệu đến các kết quả và biểu đồ trong báo cáo. Mã nguồn và dữ liệu nằm trong kho GitHub của nhóm tại \url{https://github.com/onhangocnhu/data-warehouse-assignment}; bộ dữ liệu gốc cũng có thể tải từ Kaggle.

\subsection{Yêu cầu phần mềm}

Máy cần có Docker (để chạy PostgreSQL 16), Python 3.11 trở lên (đã chạy thử với 3.14) và một bản phân phối \LaTeX{} có \code{pdflatex} với gói \code{vntex} nếu muốn biên dịch lại báo cáo và các biểu đồ. Các thư viện Python được liệt kê trong \code{src/requirements.txt}.

\subsection{Các bước chạy}

Tất cả lệnh dưới đây chạy từ thư mục gốc của kho mã nguồn. Bước 1 tạo tập tin cấu hình và khởi động cơ sở dữ liệu; nếu máy đã có sẵn một PostgreSQL chiếm cổng 5432, chỉ cần đổi \code{DB\_PORT} trong \code{.env} thành 5433.

\begin{lstlisting}[style=SQLStyle, language=bash]
# 1. Cấu hình và khởi động PostgreSQL 16 trong Docker
copy src\.env.example src\.env            # sửa DB_PASSWORD
docker compose -f src/docker-compose.yml up -d

# 2. Môi trường Python
python -m venv .venv
.venv\Scripts\activate            # Linux/macOS: source .venv/bin/activate
pip install -r src/requirements.txt

# 3. Tạo schema, bảng và chạy ETL đầy đủ (khoảng 10 giây)
python -m src.etl.run_etl --init
python -m src.etl.verify_rollback # (tùy chọn) kiểm tra tính toàn vẹn khi lỗi

# 4. EDA: khám phá dữ liệu thô và dữ liệu sau tiền xử lý
python -m src.eda.run_eda

# 5. OLAP: chạy 14 truy vấn, lưu kết quả và thời gian thực thi
python -m src.olap.run_olap

# 6. DSS: chọn mô hình, đánh giá walk-forward, chấm điểm đơn bàn giao
python -m src.dss.evaluate_models # thêm --reuse để dùng lại best_config.json
python -m src.dss.score_delay     # chạy sau bước trên

# 7. Sinh lại biểu đồ cho báo cáo (PNG 300 dpi)
python -m src.report.make_figures

# 8. Dashboard: tạo tài khoản dashboard_app (một lần) rồi mở http://localhost:8501
python -m src.dashboard.setup_role   # in 2 dòng DASHBOARD_DB_*, thêm vào .env
streamlit run src/dashboard/app.py
\end{lstlisting}

Bước 6 chạy lâu nhất (khoảng một giờ trên máy 20 luồng; \code{score\_delay} thêm khoảng 3 phút) vì phải thử 124 cấu hình trên năm tháng validation rồi huấn luyện lại mô hình cho từng tháng trong 14 tháng đánh giá; tham số được chọn được lưu ở \code{src/dss/results/best\_config.json}, nên lần chạy sau có thể bỏ qua bước tìm kiếm bằng tùy chọn \code{--reuse}. Toàn bộ chạy trên CPU, dùng dưới 2 GB bộ nhớ, không cần GPU. Có thể kiểm tra kết quả từng bước như sau: sau bước 3, \code{dm.fact\_sales} có 112.650 dòng và cả 20 tiêu chí trong \code{meta.dq\_result} đều đạt; sau bước 6, \code{dss.delay\_risk\_score} có 18.753 dòng.

Trên Linux hoặc macOS, có thể chạy dashboard bằng một lệnh \code{bash src/dashboard/run.sh}. Kịch bản này tự làm các bước 1, 2, 3, bước \code{score\_delay} (dùng \code{best\_config.json} có sẵn, không chạy lại bước tìm kiếm) và bước 8 nếu chúng chưa được làm, rồi mở dashboard; tùy chọn \code{--check} chỉ chuẩn bị và kiểm tra mà không mở dashboard.

Trên Windows bật Smart App Control, thư viện \code{psycopg} có thể bị chặn (lỗi \textit{DLL load failed\ldots Application Control policy}). Khi đó đặt biến môi trường \code{PSYCOPG\_IMPL=python} và thêm thư mục \code{bin} của một bản PostgreSQL cài trên máy vào \code{PATH} để dùng \code{libpq.dll} của bản đó.

\subsection{Cấu trúc mã nguồn}

\begin{lstlisting}[style=SQLStyle, language={}]
data/raw/                 9 tập tin CSV gốc của Olist
data/reference/           bảng vùng của 27 bang, bản dịch 2 ngành hàng còn thiếu
src/sql/                  00-05 tạo schema và bảng; 10 transform; 20 nạp bảng chiều;
                          30 nạp bảng fact; 40 kiểm tra chất lượng; 50 truy vấn OLAP
src/etl/                  run_etl.py (điều phối), extract.py (COPY), db.py (nhật ký),
                          verify_rollback.py
src/eda/                  run_eda.py; results/ chứa kết quả khám phá dữ liệu
src/olap/                 run_olap.py; results/ chứa kết quả và thời gian từng truy vấn
src/dss/                  features.py (đặc trưng point-in-time), evaluate_models.py,
                          score_delay.py; results/ chứa kết quả đánh giá
src/dashboard/            app.py (điều hướng), db.py, filters.py, ui.py, views/ (7 trang),
                          setup_role.py, run.sh (chạy dashboard bằng một lệnh)
src/report/               make_figures.py: sinh biểu đồ của báo cáo từ cơ sở dữ liệu
report/                   mã nguồn LaTeX của báo cáo; Sections/hinh/ chứa các sơ đồ TikZ
src/docker-compose.yml    PostgreSQL 16
\end{lstlisting}

\section{Phụ lục: Bảng mã bang của Brazil}\label{sec:mabang}

Bảng~\ref{tab:mabang} liệt kê 27 đơn vị hành chính cấp bang của Brazil (26 bang và Quận Liên bang DF) theo mã hai chữ cái dùng trong dữ liệu và trong báo cáo, cùng vùng mà nhóm gán cho từng bang trong bảng chiều \code{dim\_state}.

\begin{table}[H]
\centering
\small
\caption{Mã, tên và vùng của 27 bang Brazil}\label{tab:mabang}
\begin{tabular}{|c|l|l||c|l|l|}
\hline
\textbf{Mã} & \textbf{Tên bang} & \textbf{Vùng} & \textbf{Mã} & \textbf{Tên bang} & \textbf{Vùng}\\ \hline
AC & Acre & Norte & PB & Paraíba & Nordeste\\ \hline
AL & Alagoas & Nordeste & PE & Pernambuco & Nordeste\\ \hline
AM & Amazonas & Norte & PI & Piauí & Nordeste\\ \hline
AP & Amapá & Norte & PR & Paraná & Sul\\ \hline
BA & Bahia & Nordeste & RJ & Rio de Janeiro & Sudeste\\ \hline
CE & Ceará & Nordeste & RN & Rio Grande do Norte & Nordeste\\ \hline
DF & Distrito Federal & Centro-Oeste & RO & Rondônia & Norte\\ \hline
ES & Espírito Santo & Sudeste & RR & Roraima & Norte\\ \hline
GO & Goiás & Centro-Oeste & RS & Rio Grande do Sul & Sul\\ \hline
MA & Maranhão & Nordeste & SC & Santa Catarina & Sul\\ \hline
MG & Minas Gerais & Sudeste & SE & Sergipe & Nordeste\\ \hline
MS & Mato Grosso do Sul & Centro-Oeste & SP & São Paulo & Sudeste\\ \hline
MT & Mato Grosso & Centro-Oeste & TO & Tocantins & Norte\\ \hline
PA & Pará & Norte & & & \\ \hline
\end{tabular}
\end{table}
